% GENERATED FILE. Edit canonical lessons, then run npm run generate:books.
% canonical-source-hash: fnv1a64:9a1726c92618c52c
% canonical-lessons: MW-C100-darvajo, MW-C100-pol, MW-C100-haveli, MW-C100-melo, MW-C100-hat

\chapter{Gate, Gateway, Mansion, Fair, Weekly market}
\label{ch:mw-gate-gateway-mansion-fair-weekly-market}
\hlchaptermodality{100}

\begin{chapteropening}
\textbf{By the end of this chapter:} \emph{I can say a gate, a gateway, a mansion, a fair and a weekly market.}

Complete the last lesson of chapter 100: I can say a gate, a gateway, a mansion, a fair and a weekly market.
\end{chapteropening}

\section[\texorpdfstring{\mw{दरवाजो} (darvājo)}{दरवाजो (darvājo)}]{\mw{दरवाजो} (darvājo) — a gate, a doorway}

\label{lesson:MW-C100-darvajo}



\begin{quote}
Before the new one: say the Marwadi for a stepwell, then the Marwadi for a pond, a lake.
\end{quote}

\subsection*{You'll want to know: \mw{दरवाजो}}
\textbf{\mw{दरवाजो}} — \emph{darvājo} — \textquotedblleft{}a gate, a doorway\textquotedblright{}.

\begin{cousinweb}
\textbf{\mw{दरवाजो}} (\emph{darvājo}) is a loan, usually traced to Persian \emph{darvāza}, \textquotedblleft{}gate, door\textquotedblright{}. Marwadi says the \emph{z} as \emph{j}. The Persian word is a cousin of English \emph{door}.
\end{cousinweb}

\begin{grammarlens}[title={What you've built}]
One more word for this chapter.
\end{grammarlens}

\subsection*{Your turn}
\begin{itemize}
\raggedright
  \item \emph{Say it:} \emph{darvājo}
  \item \emph{Say it:} \emph{darvājo}, once more
  \item \emph{Say it:} say \emph{darvājo}
  \item \emph{Recall:} say the Marwadi for a stepwell, then the Marwadi for a pond, a lake, then say \emph{darvājo} again
\end{itemize}

\begin{tcolorbox}[breakable,colback=teal!4,colframe=teal!35!black,title={Before you move on}]
What does \mw{दरवाजो} mean? (\textquotedblleft{}A gate, a doorway\textquotedblright{}.) Say it once more.
\end{tcolorbox}
\section[\texorpdfstring{\mw{पोळ} (poḷ)}{पोळ (poḷ)}]{\mw{पोळ} (poḷ) — a gateway}

\label{lesson:MW-C100-pol}



\begin{quote}
Before the new one: say the Marwadi for a pond, a lake, then the Marwadi for a gate.
\end{quote}

\subsection*{You'll want to know: \mw{पोळ}}
\textbf{\mw{पोळ}} — \emph{poḷ} — \textquotedblleft{}a gateway\textquotedblright{}.

\begin{cousinweb}
\textbf{\mw{पोळ}} (\emph{poḷ}) is a tadbhava, a word inherited from Sanskrit and reshaped by sound change. It comes from \emph{pratolī}, \textquotedblleft{}gate of town or fort\textquotedblright{}. CDIAL gives no Marwari form in this entry (CDIAL 8633); the match is through Hindi \emph{pol} and Gujarati \emph{poḷ}.
\end{cousinweb}

\begin{grammarlens}[title={What you've built}]
One more word for this chapter.
\end{grammarlens}

\subsection*{Your turn}
\begin{itemize}
\raggedright
  \item \emph{Say it:} \emph{poḷ}
  \item \emph{Say it:} \emph{poḷ}, once more
  \item \emph{Say it:} say \emph{poḷ}
  \item \emph{Recall:} say the Marwadi for a pond, a lake, then the Marwadi for a gate, then say \emph{poḷ} again
\end{itemize}

\begin{tcolorbox}[breakable,colback=teal!4,colframe=teal!35!black,title={Before you move on}]
What does \mw{पोळ} mean? (\textquotedblleft{}A gateway\textquotedblright{}.) Say it once more.
\end{tcolorbox}
\section[\texorpdfstring{\mw{हवेली} (havelī)}{हवेली (havelī)}]{\mw{हवेली} (havelī) — a mansion}

\label{lesson:MW-C100-haveli}



\begin{quote}
Before the new one: say the Marwadi for a gate, then the Marwadi for a gateway.
\end{quote}

\subsection*{You'll want to know: \mw{हवेली}}
\textbf{\mw{हवेली}} — \emph{havelī} — \textquotedblleft{}a mansion\textquotedblright{}.

\begin{cousinweb}
\textbf{\mw{हवेली}} (\emph{havelī}) is a loan, usually traced to Persian \emph{havēlī}, \textquotedblleft{}mansion\textquotedblright{}. Persian took it from Arabic \emph{ḥawālī}, \textquotedblleft{}surroundings\textquotedblright{} (Lane, root \emph{ḥ-w-l}).
\end{cousinweb}

\begin{grammarlens}[title={What you've built}]
One more word for this chapter.
\end{grammarlens}

\subsection*{Your turn}
\begin{itemize}
\raggedright
  \item \emph{Say it:} \emph{havelī}
  \item \emph{Say it:} \emph{havelī}, once more
  \item \emph{Say it:} say \emph{havelī}
  \item \emph{Recall:} say the Marwadi for a gate, then the Marwadi for a gateway, then say \emph{havelī} again
\end{itemize}

\begin{tcolorbox}[breakable,colback=teal!4,colframe=teal!35!black,title={Before you move on}]
What does \mw{हवेली} mean? (\textquotedblleft{}A mansion\textquotedblright{}.) Say it once more.
\end{tcolorbox}
\section[\texorpdfstring{\mw{मेळो} (meḷo)}{मेळो (meḷo)}]{\mw{मेळो} (meḷo) — a fair}

\label{lesson:MW-C100-melo}



\begin{quote}
Before the new one: say the Marwadi for a gateway, then the Marwadi for a mansion.
\end{quote}

\subsection*{You'll want to know: \mw{मेळो}}
\textbf{\mw{मेळो}} — \emph{meḷo} — \textquotedblleft{}a fair\textquotedblright{}.

\begin{cousinweb}
\textbf{\mw{मेळो}} (\emph{meḷo}) is a tadbhava, a word inherited from Sanskrit and reshaped by sound change. It comes from \emph{mela}, \textquotedblleft{}meeting\textquotedblright{}. CDIAL lists Marwari \emph{meḷo} under it (CDIAL 10331). It is built on the root \emph{mil}, \textquotedblleft{}to meet\textquotedblright{}.
\end{cousinweb}

\begin{grammarlens}[title={What you've built}]
One more word for this chapter.
\end{grammarlens}

\subsection*{Your turn}
\begin{itemize}
\raggedright
  \item \emph{Say it:} \emph{meḷo}
  \item \emph{Say it:} \emph{meḷo}, once more
  \item \emph{Say it:} say \emph{meḷo}
  \item \emph{Recall:} say the Marwadi for a gateway, then the Marwadi for a mansion, then say \emph{meḷo} again
\end{itemize}

\begin{tcolorbox}[breakable,colback=teal!4,colframe=teal!35!black,title={Before you move on}]
What does \mw{मेळो} mean? (\textquotedblleft{}A fair\textquotedblright{}.) Say it once more.
\end{tcolorbox}
\section[\texorpdfstring{\mw{हाट} (hāṭ)}{हाट (hāṭ)}]{\mw{हाट} (hāṭ) — a weekly market}

\label{lesson:MW-C100-hat}



\begin{quote}
Before the new one: say the Marwadi for a mansion, then the Marwadi for a fair.
\end{quote}

\subsection*{You'll want to know: \mw{हाट}}
\textbf{\mw{हाट}} — \emph{hāṭ} — \textquotedblleft{}a weekly market\textquotedblright{}.

\begin{cousinweb}
\textbf{\mw{हाट}} (\emph{hāṭ}) is a tadbhava, a word inherited from Sanskrit and reshaped by sound change. It comes from \emph{haṭṭa}, \textquotedblleft{}market, fair\textquotedblright{}. CDIAL lists Old Marwari \emph{hāṭi} under it (CDIAL 13944).
\end{cousinweb}

\begin{grammarlens}[title={What you've built}]
One more word for this chapter.
\end{grammarlens}

\subsection*{Your turn}
\begin{itemize}
\raggedright
  \item \emph{Say it:} \emph{hāṭ}
  \item \emph{Say it:} \emph{hāṭ}, once more
  \item \emph{Say it:} say \emph{hāṭ}
  \item \emph{Recall:} say the Marwadi for a mansion, then the Marwadi for a fair, then say \emph{hāṭ} again
\end{itemize}

\begin{tcolorbox}[breakable,colback=teal!4,colframe=teal!35!black,title={Before you move on}]
What does \mw{हाट} mean? (\textquotedblleft{}A weekly market\textquotedblright{}.) Say it once more.
\end{tcolorbox}
